% GENERATED FILE. Edit canonical lessons, then run npm run generate:books.
% canonical-source-hash: fnv1a64:f204943b1d802ca1
% canonical-lessons: ZH-C147-houtian, ZH-W147-hou-after, ZH-C147-hui-future, ZH-C147-keneng, ZH-C147-yiding, ZH-C147-tianqi-yubao, ZH-C147-practice

\chapter{It Might Rain}
\label{ch:zh-it-might-rain}
\hlchaptermodality{147}

\begin{chapteropening}
\textbf{By the end of this chapter:} \emph{I can say what will probably happen, and what the weather will do.}

Saying what will happen with \zh{会}, maybe and certainly, the day after tomorrow and the weather forecast, built on one new written character: \zh{后}.
\end{chapteropening}

\section[\texorpdfstring{\zh{后天} (hòutiān)}{后天 (hòutiān)}]{\zh{后天} (hòutiān) — the day after tomorrow}

\label{lesson:ZH-C147-houtian}



\begin{quote}
A new chapter, and it opens with a day that has not come yet.

\begin{itemize}
\raggedright
  \item \emph{Recall:} say \emph{píngshí}, usually — \textbf{R1}, the chapter you have finished
  \item \emph{Recall:} say \zh{明天} — \emph{míngtiān}, tomorrow — \textbf{R3}, durable
  \item \emph{Recall:} say \emph{hòu}, after; behind — \textbf{R3}
\end{itemize}
\end{quote}

\subsection*{You'll want to know: \zh{后天}}
\textbf{\zh{后天}} — \emph{hòutiān} — \textquotedblleft{}the day after tomorrow\textquotedblright{}.

\begin{quote}
\textbf{\zh{后天我们去}} \emph{hòutiān wǒmen qù} — We're going the day after tomorrow.
\end{quote}

\begin{quote}
\textbf{\zh{今天} \zh{明天} \zh{后天}} \emph{jīntiān, míngtiān, hòutiān} — today, tomorrow, the day after
\end{quote}

The next lesson puts the first character in your hand. It is the \emph{hòu} you already say for behind and after.

\begin{cousinweb}
\textbf{\zh{后天}} (\emph{hòutiān}) is built from \zh{后} 'after' and \zh{天} 'day'. Both are inherited from Middle Chinese, \emph{huwX} and \emph{then} (Guangyun, in Baxter's transcription). The Guangyun lists \emph{huwX} under the traditional form of the character and glosses it 'before and after'; it glosses \emph{then} 'the heavens above', so 'day' is assumed here, not quoted.
\end{cousinweb}

\begin{grammarlens}[title={What you've built}]
\zh{后天} takes its place in the row of days you already say: \emph{zuótiān}, \zh{今天}, \zh{明天}, \zh{后天}.

Like the others it goes early in the sentence, before the verb: \textbf{\zh{后天我们去}} \emph{hòutiān wǒmen qù}.
\end{grammarlens}

\subsection*{Your turn}
\begin{itemize}
\raggedright
  \item \emph{Say it:} \emph{hòutiān}
  \item \emph{Say it:} \emph{hòutiān}, once more
  \item \emph{Say it:} \zh{后天我们去} — \emph{hòutiān wǒmen qù} — We're going the day after tomorrow.
  \item \emph{Say it:} \zh{今天} \zh{明天} \zh{后天} — \emph{jīntiān, míngtiān, hòutiān} — today, tomorrow, the day after
  \item \emph{Recall:} say \emph{píngshí}, then \emph{hòutiān} again
\end{itemize}

\begin{tcolorbox}[breakable,colback=teal!4,colframe=teal!35!black,title={Before you move on}]
What does \zh{后天} mean? (\textquotedblleft{}The day after tomorrow\textquotedblright{}.) And the day before it? (\textbf{\zh{明天}}.)
\end{tcolorbox}
\section[\texorpdfstring{\zh{后} (hòu)}{后 (hòu)}]{\zh{后} — a three-stroke frame, then \zh{口}}

\label{lesson:ZH-W147-hou-after}



\begin{quote}
The word first, then the character inside it.

\begin{itemize}
\raggedright
  \item \emph{Recall:} say \emph{hòutiān} — \textbf{R1}, one lesson back, before it decays
\end{itemize}

\begin{quote}
\zh{口} \zh{天}
\end{quote}

\emph{Read it:} each one aloud, with nothing said first — \emph{kǒu}, \emph{tiān}
\end{quote}

\subsection*{Script — a three-stroke frame, then \zh{口}}
\hlblockfigure{\detokenize{figures/ZH-W147-hou-after-filmstrip.pdf}}{How \zh{后} is written, stroke by stroke}

\begin{quote}
\zh{后}
\end{quote}

\textbf{Two pieces, and the second you already write.}

\textbf{The frame, three strokes}: a short left-falling stroke along the top, from right to left; a long left-falling stroke down the left side, curving out to the lower left; then a horizontal across the middle, from the long stroke to the right.

\textbf{\zh{口}, three strokes}, under that horizontal, on the right: the left side, the top that turns down the right side, the bottom bar.

3 + 3 = six.

\subsection*{Your turn}
\begin{itemize}
\raggedright
  \item \emph{Write it:} \zh{后}, all six, with the model in front of you
  \item \emph{Cover:} write it again from memory, and count the pieces, not the strokes
\end{itemize}

\begin{tcolorbox}[breakable,colback=teal!4,colframe=teal!35!black,title={Before you move on}]
How many strokes in \zh{后}? (\textbf{Six}.) Which stroke comes first? (\textbf{The short one along the top}.) Which character you write sits inside the frame? (\textbf{\zh{口}}.)

Source: \href{https://raw.githubusercontent.com/chanind/hanzi-writer-data/68d10a4b21150cae5e1ebbd223eed289cf32d90c/data/\%E5\%90\%8E.json}{Hanzi Writer Data \zh{后}, pinned stroke paths and medians}.
\end{tcolorbox}
\section[\texorpdfstring{\zh{明天会下雨} (míngtiān huì xiàyǔ)}{明天会下雨 (míngtiān huì xiàyǔ)}]{\zh{明天会下雨} — \zh{会} for what will happen}

\label{lesson:ZH-C147-hui-future}



\begin{quote}
A word you already know does a second job in this lesson.

\begin{itemize}
\raggedright
  \item \emph{Read it:} \zh{后天}, aloud — \emph{hòutiān}
  \item \emph{Recall:} say \zh{我会开车} — \emph{wǒ huì kāichē}, I can drive
  \item \emph{Recall:} say \zh{下雨} — \emph{xiàyǔ}, to rain
\end{itemize}
\end{quote}

\begin{grammarlens}[title={Grammar Lens: \zh{会} before a verb, for what is going to happen}]
\begin{quote}
\textbf{\zh{会} before a verb can say what will probably happen, or what someone will do.}
\end{quote}

\begin{quote}
\textbf{\zh{明天会下雨}} \emph{míngtiān huì xiàyǔ} — It'll rain tomorrow. \textbf{\zh{后天不会下雨}} \emph{hòutiān bú huì xiàyǔ} — It won't rain the day after tomorrow.
\end{quote}

\zh{不} before the fall of \zh{会} is \emph{bú}, as in \emph{bú qù}.

You met \zh{会} as \emph{can}, for something learned: \zh{我会开车}. It is the same character in a different use. The rain has learned nothing, so \zh{明天会下雨} is a forecast, not a skill; a time word in front usually tells you which one you hear.

You have said the future without \zh{会} for a long time: \textbf{\zh{明天我们去}} \emph{míngtiān wǒmen qù} needs no extra word. \zh{会} adds a forecast, what you expect to happen:

\begin{quote}
\textbf{\zh{明天我会去}} \emph{míngtiān wǒ huì qù} — I'll go tomorrow.
\end{quote}
\end{grammarlens}

\subsection*{Your turn}
Say these aloud:

\begin{itemize}
\raggedright
  \item \zh{明天会下雨} — \emph{míngtiān huì xiàyǔ} — It'll rain tomorrow.
  \item \zh{后天不会下雨} — \emph{hòutiān bú huì xiàyǔ} — It won't rain the day after tomorrow.
  \item \zh{明天我会去} — \emph{míngtiān wǒ huì qù} — I'll go tomorrow.
  \item \zh{我会开车} — \emph{wǒ huì kāichē} — I can drive: the other \zh{会}.
\end{itemize}

\begin{tcolorbox}[breakable,colback=teal!4,colframe=teal!35!black,title={Before you move on}]
What does \zh{会} add to \zh{明天会下雨}? (\textbf{A forecast: it will rain}.) How do you say it won't? (\textbf{\zh{不会}}, said \emph{bú huì}.) What else can \zh{会} mean? (\textbf{Can, for a skill: \zh{我会开车}}.)
\end{tcolorbox}
\section[kěnéng]{kěnéng — maybe; possibly}

\label{lesson:ZH-C147-keneng}



\begin{quote}
Before the new one: read the last lesson's sentence, then a word that shares a sound with this one.

\begin{itemize}
\raggedright
  \item \emph{Read it:} \zh{明天会下雨}, aloud — \emph{míngtiān huì xiàyǔ}
  \item \emph{Recall:} say \emph{kěyǐ}, may — \textbf{R3}, durable
\end{itemize}
\end{quote}

\subsection*{You'll want to know: kěnéng}
\emph{kěnéng} — \textquotedblleft{}maybe; possibly\textquotedblright{}.

\begin{quote}
\emph{tā kěnéng bù lái} — He may not come.
\end{quote}

\begin{quote}
\emph{míngtiān kěnéng huì xiàyǔ} — It might rain tomorrow.
\end{quote}

Its first character comes later; for now it is a word for the ear.

\begin{cousinweb}
\emph{kěnéng} is built from kě 'may', the kě of \emph{kěyǐ}, and \zh{能} 'can', the \zh{能} you write. Both are inherited from Middle Chinese, \emph{khaX} and \emph{nong} (Guangyun, in Baxter's transcription), glossed 'to permit' and 'skilled'. 'Maybe' is the pair's own sense, which neither gloss gives.
\end{cousinweb}

\begin{grammarlens}[title={What you've built}]
\emph{kěnéng} goes before the verb, or before \zh{会}: \emph{kěnéng huì xiàyǔ}, it might rain.

With \zh{会} alone the forecast sounds sure; with \emph{kěnéng} in front it is only a chance.
\end{grammarlens}

\subsection*{Your turn}
\begin{itemize}
\raggedright
  \item \emph{Say it:} \emph{kěnéng}
  \item \emph{Say it:} \emph{kěnéng}, once more
  \item \emph{Say it:} \emph{tā kěnéng bù lái} — He may not come.
  \item \emph{Say it:} \emph{míngtiān kěnéng huì xiàyǔ} — It might rain tomorrow.
  \item \emph{Recall:} say \emph{kěyǐ}, then \emph{kěnéng}, and hear the same kě
\end{itemize}

\begin{tcolorbox}[breakable,colback=teal!4,colframe=teal!35!black,title={Before you move on}]
What does \emph{kěnéng} mean? (\textquotedblleft{}Maybe; possibly\textquotedblright{}.) Where does it go? (\textbf{Before the verb, or before \zh{会}}.)
\end{tcolorbox}
\section[yídìng]{yídìng — certainly; definitely}

\label{lesson:ZH-C147-yiding}



\begin{quote}
Before the new one: last lesson's word, then \zh{一} with its rising tone.

\begin{itemize}
\raggedright
  \item \emph{Recall:} say \emph{kěnéng}, maybe — \textbf{R1}, one lesson back
  \item \emph{Recall:} say \zh{一个} — \emph{yí ge}, one of them, with \zh{一} rising — \textbf{R3}, durable
\end{itemize}
\end{quote}

\subsection*{You'll want to know: yídìng}
\emph{yídìng} — \textquotedblleft{}certainly; definitely\textquotedblright{}.

\begin{quote}
\emph{tā yídìng huì lái} — He'll definitely come.
\end{quote}

\zh{一} rises before the fall of dìng: \emph{yídìng}, as in \emph{yí ge}.

Its second character comes later; for now it is a word for the ear.

\begin{cousinweb}
\emph{yídìng} is built from \zh{一} 'one' and dìng 'settled'. Both are inherited from Middle Chinese, \emph{'jit} and \emph{dengH} (Guangyun, in Baxter's transcription), \emph{dengH} glossed 'calm; settled'. 'Certainly' is the pair's own sense, not the gloss's.
\end{cousinweb}

\begin{grammarlens}[title={What you've built}]
\emph{kěnéng} and \emph{yídìng} take the same place, before \zh{会} or the verb: a chance, then a certainty.

\begin{itemize}
\raggedright
  \item \emph{tā kěnéng huì lái} — he may come.
  \item \emph{tā yídìng huì lái} — he will, for sure.
\end{itemize}
\end{grammarlens}

\subsection*{Your turn}
\begin{itemize}
\raggedright
  \item \emph{Say it:} \emph{yídìng}
  \item \emph{Say it:} \emph{yídìng}, once more
  \item \emph{Say it:} \emph{tā yídìng huì lái} — He'll definitely come.
  \item \emph{Say it:} \emph{wǒ yídìng qù} — I'll be there, for sure.
  \item \emph{Recall:} say \emph{kěnéng}, then \emph{yídìng}: a chance, then a certainty
\end{itemize}

\begin{tcolorbox}[breakable,colback=teal!4,colframe=teal!35!black,title={Before you move on}]
What does \emph{yídìng} mean? (\textquotedblleft{}Certainly; definitely\textquotedblright{}.) What does \zh{一} do before the fall of dìng? (\textbf{It rises: \emph{yí}}.)
\end{tcolorbox}
\section[tiānqì yùbào]{tiānqì yùbào — weather forecast}

\label{lesson:ZH-C147-tianqi-yubao}



\begin{quote}
Before the new one: last lesson's word, then the weather itself.

\begin{itemize}
\raggedright
  \item \emph{Recall:} say \emph{yídìng}, certainly — \textbf{R1}, one lesson back
  \item \emph{Recall:} say \zh{天气} — \emph{tiānqì}, the weather — \textbf{R3}, durable
\end{itemize}
\end{quote}

\subsection*{You'll want to know: tiānqì yùbào}
\emph{tiānqì yùbào} — \textquotedblleft{}weather forecast\textquotedblright{}.

\begin{quote}
\emph{tiānqì yùbào shuō míngtiān huì xiàyǔ} — The forecast says it'll rain tomorrow.
\end{quote}

Its last two characters come later; for now it is a word for the ear.

\begin{cousinweb}
\emph{tiānqì yùbào} is built from \zh{天气} 'weather' and yùbào. yùbào is built from yù and bào, inherited from Middle Chinese \emph{yoH} and \emph{pawH} (Guangyun, in Baxter's transcription), glossed 'beforehand' and 'to report'. A report given beforehand is the plain reading of the two; 'forecast' for the whole is the modern compound's own sense.
\end{cousinweb}

\begin{grammarlens}[title={What you've built}]
\emph{tiānqì yùbào shuō …}, the forecast says, and then the forecast itself, with \zh{会}:

\begin{quote}
\emph{tiānqì yùbào shuō hòutiān bú huì xiàyǔ} — The forecast says it won't rain the day after tomorrow.
\end{quote}
\end{grammarlens}

\subsection*{Your turn}
\begin{itemize}
\raggedright
  \item \emph{Say it:} \emph{tiānqì yùbào}
  \item \emph{Say it:} \emph{tiānqì yùbào}, once more
  \item \emph{Say it:} \emph{tiānqì yùbào shuō míngtiān huì xiàyǔ} — The forecast says it'll rain tomorrow.
  \item \emph{Recall:} say \zh{天气}, then \emph{tiānqì yùbào}
\end{itemize}

\begin{tcolorbox}[breakable,colback=teal!4,colframe=teal!35!black,title={Before you move on}]
What does \emph{tiānqì yùbào} mean? (\textquotedblleft{}Weather forecast\textquotedblright{}.) What does the forecast say, with \zh{会}? (\emph{tiānqì yùbào shuō míngtiān huì xiàyǔ}.)
\end{tcolorbox}
\section[Practice]{\zh{明天会下雨吗} — \zh{后天不会}}

\label{lesson:ZH-C147-practice}



\begin{quote}
The chapter closes here, with two people planning around the rain.

\begin{itemize}
\raggedright
  \item \emph{Recall:} say \emph{tiānqì yùbào}, the weather forecast — \textbf{R1}, one lesson back
  \item \emph{Recall:} say \emph{yídìng}, certainly — \textbf{R1}
  \item \emph{Recall:} say \emph{kěnéng}, maybe — \textbf{R2}
\end{itemize}
\end{quote}

\subsection*{The exchange}
\begin{quote}
\textbf{A:} \zh{明天会下雨吗}
\end{quote}

\begin{quote}
\emph{míngtiān huì xiàyǔ ma} — Will it rain tomorrow?
\end{quote}

\begin{quote}
\textbf{B:} \emph{huì, tiānqì yùbào shuō huì xiàyǔ} — It will. The forecast says so.
\end{quote}

\begin{quote}
\textbf{A:} \zh{后天呢}
\end{quote}

\begin{quote}
\emph{hòutiān ne} — And the day after?
\end{quote}

\begin{quote}
\textbf{B:} \zh{后天不会}
\end{quote}

\begin{quote}
\emph{hòutiān bú huì} — Not the day after.
\end{quote}

\begin{quote}
\textbf{A:} \zh{我们后天去}
\end{quote}

\begin{quote}
\emph{wǒmen hòutiān qù} — Then we'll go the day after tomorrow.
\end{quote}

\begin{quote}
\textbf{B:} \zh{他会来吗}
\end{quote}

\begin{quote}
\emph{tā huì lái ma} — Will he come?
\end{quote}

\begin{quote}
\textbf{A:} \emph{tā kěnéng bù lái} — He may not.
\end{quote}

\begin{quote}
\textbf{B:} \emph{wǒ yídìng lái} — I'll come, for sure.
\end{quote}

\zh{会} answers a question about \zh{会} on its own: \emph{míngtiān huì xiàyǔ ma?} — \emph{huì}. The answer drops everything else, as \zh{是} and \zh{有} do.

\begin{tcolorbox}[breakable,skin=enhanced,colback=orange!6,colframe=orange!45!black,boxrule=0.5pt,arc=1mm,left=6pt,right=6pt,top=4pt,bottom=4pt,fonttitle=\bfseries,title={Read this}]
\emph{Read it:} the exchange through once without stopping, then answer the question below
\end{tcolorbox}

\begin{grammarlens}[title={What you've built — a forecast, a chance, a certainty}]
Three ways to talk about what is coming, all before the verb:

\begin{itemize}
\raggedright
  \item \textbf{\zh{明天会下雨}} \emph{míngtiān huì xiàyǔ} — it will rain tomorrow.
  \item \emph{míngtiān kěnéng huì xiàyǔ} — it might.
  \item \emph{tā yídìng huì lái} — he certainly will come.
\end{itemize}

\textbf{Listen.} \emph{Hòutiān huì xiàyǔ ma?}

\textbf{Speak.} Say what the weather will do tomorrow, then what you will do the day after.

\emph{Write it:} \zh{后天} — the frame first, then \zh{口}, then \zh{天}
\end{grammarlens}

\subsection*{Your turn}
\begin{itemize}
\raggedright
  \item \emph{Say it:} \zh{明天会下雨吗} — \emph{míngtiān huì xiàyǔ ma} — Will it rain tomorrow?
  \item \emph{Say it:} \zh{后天不会} — \emph{hòutiān bú huì} — Not the day after.
  \item \emph{Say it:} the whole exchange, both parts
  \item \emph{Write it:} \zh{明天会下雨}
\end{itemize}

\begin{tcolorbox}[breakable,colback=teal!4,colframe=teal!35!black,title={Before you move on}]
Say it will rain tomorrow. (\textbf{\zh{明天会下雨}}.) Say it won't rain the day after. (\textbf{\zh{后天不会下雨}}.) Say he may not come. (\emph{tā kěnéng bù lái}.)
\end{tcolorbox}
